% ============================================================
\chapter{Experimental Protocols and Evaluation Methods}
\label{ch:protocols}
% ============================================================

\section{Experimental Phases}

The complete experimental campaign is separated into seven phases so that an
exploratory observation is not mistaken for controlled evidence.

\begin{figure}[H]
\centering
\resizebox{\textwidth}{!}{\begin{tikzpicture}[
  node distance=0.58cm,
  phase/.style={draw,rounded corners=2pt,text width=2.15cm,minimum height=1.2cm,
    align=center,font=\small,inner sep=4pt},
  control/.style={draw,dashed,rounded corners=2pt,text width=8.6cm,minimum height=0.8cm,
    align=center,font=\footnotesize,inner sep=3pt},
  arrow/.style={-{Latex[length=2.4mm]},thick,draw=azulUC3M}
]
\node[phase,fill=softblue] (explore) {\textbf{Phase 1}\\PIDNet exploratory development};
\node[phase,fill=softgreen,right=of explore] (ablate) {\textbf{Phase 2}\\single-factor PIDNet ablations};
\node[phase,fill=softorange,right=of ablate] (freeze) {\textbf{Phase 3}\\checkpoint audit and reproduction};
\node[phase,fill=softpurple,right=of freeze] (rod) {\textbf{Phase 4}\\PIDNet--ROD repeatability};
\node[phase,fill=softblue,below=1.1cm of rod] (wide) {\textbf{Phase 5}\\nine architectures, two policies};
\node[phase,fill=softorange,left=of wide] (converge) {\textbf{Phase 6}\\blue+green convergence follow-up};
\node[phase,fill=softgreen,left=of converge] (deploy) {\textbf{Phase 7}\\CPU, H100, TensorRT};
\draw[arrow] (explore) -- (ablate);
\draw[arrow] (ablate) -- (freeze);
\draw[arrow] (freeze) -- (rod);
\draw[arrow] (rod) -- (wide);
\draw[arrow] (wide) -- (converge);
\draw[arrow] (converge) -- (deploy);
\node[control,below=0.55cm of converge] {
  \textbf{Controls:} Phases 1--3 retain one task and split while the PIDNet
  configuration is developed and audited. Phases 4--6 use matched splits,
  objectives, seeds, and thresholds; Phase 6 replaces the early fixed budget
  with a declared stopping rule.
};
\end{tikzpicture}
}
\caption[Experimental phases and controls.]{Experimental phases and their
principal controls. Dataset identity, target policy, and metric definitions
remain fixed as the question changes from configuration selection to model
comparison and runtime measurement.}
\label{fig:experimental_design}
\end{figure}

Figure~\ref{fig:experimental_design} provides the overview; the paragraphs
below state the purpose and controls of each phase.

\textbf{Phase 1: exploratory PIDNet development.} Five runs changed multiple
settings to obtain a working pipeline and reveal promising or problematic
directions. These runs guide questions; they do not isolate causal factors.

\textbf{Phase 2: controlled PIDNet ablation.} A reference configuration and
four one-factor variants test augmentation, false-safe penalty, resolution,
and class weights. All other recorded settings are held fixed. One seed is
used, so this is a controlled engineering comparison without a variance
estimate.

\textbf{Phase 3: checkpoint audit.} The validation-selected augmentation-off
PIDNet checkpoint is identified, reevaluated, qualitatively ranked, and
benchmarked. Subsequent measurements refer to these saved parameters.

\textbf{Phase 4: PIDNet--ROD repeatability.} Both architectures are
trained for seeds 1337, 2027, and 4242 under blue-only and blue+green policies
with one controlled recipe. Separate ROD runs use the optimisation settings
that can be reconstructed from the publication.

\textbf{Phase 5: architecture matrix.} Seven additional systems complete a
nine-model, three-seed, two-policy matrix. Every policy contains 27 scored
model--seed records; together the two ledgers contain 54.

\textbf{Phase 6: convergence-aware follow-up.} The nine-model, three-seed
matrix is repeated for the principal blue+green policy with a 300-epoch
ceiling, 60-epoch cosine decay, a 60-epoch minimum, and validation-based early
stopping. This phase contains another 27 scored model--seed records.

\textbf{Phase 7: runtime evaluation.} Models are timed on a Ryzen~5~5500 CPU
and a common H100 MIG server protocol. The reference PIDNet checkpoint and the
reconstructed ROD checkpoint are also exported to TensorRT and evaluated on an
RTX~5060.

All experiments exclude the official CaT test images from gradient updates and
automatic checkpoint selection. The test set is nevertheless evaluated for
multiple exploratory, ablation, and architecture runs over the lifetime of the
project. It should therefore be understood as a consistently held-out
benchmark partition, not as a single-use blinded set whose scores were hidden
until every engineering decision was complete. Model checkpoints are selected
by validation mIoU, but repeated human inspection of test reports remains a
possible source of research-level selection bias.

\section{Training Procedure}

Each scored run follows the same operational sequence:
\begin{enumerate}
  \item load one complete YAML configuration and set the Python, NumPy, and
  PyTorch random seeds;
  \item construct training and validation datasets from the processed
  manifest, enabling augmentation only for training when configured;
  \item compute class weights from training masks only, or load explicit
  neutral weights;
  \item build the named two-class model and AdamW optimiser;
  \item for every epoch, shuffle the training loader, compute the configured
  loss, backpropagate, update weights, and advance the scheduler;
  \item run an unaugmented validation pass, accumulating its confusion matrix;
  \item write the epoch history and \texttt{last.pt}; replace
  \texttt{best.pt} if validation mIoU exceeds the previous maximum;
  \item after training, load \texttt{best.pt} in the evaluator and score the
  complete test split without gradients or shuffling.
\end{enumerate}

Non-finite training loss or validation metrics raise an error rather than
silently saving a checkpoint. Mixed precision uses bfloat16 when CUDA supports
it and falls back to float16 otherwise. The ordinary 15-epoch experiments do
not use early stopping; every model receives the same maximum number of
training passes. This fixed budget improves protocol matching but is not
necessarily optimal for every architecture.

The convergence-aware blue+green follow-up uses the same operational sequence
with three changes. The epoch ceiling is 300, cosine decay reaches the minimum
learning rate at epoch 60 and then holds it constant, and early stopping is
enabled only after epoch 60. The stopping counter requires an improvement of
at least $10^{-4}$ in validation mIoU and has patience 25. Checkpoint saving
uses any strict improvement, so it can record a slightly higher value that is
too small to reset the patience counter. Every reported test result still
comes from \texttt{best.pt}, never from the final in-memory epoch.

The test evaluator computes weighted cross-entropy loss using the class weights
stored in the checkpoint. Quality metrics are derived independently from
thresholded predictions, so the test loss and mIoU need not rank models in the
same order. The prediction interface is normalised to one full-resolution
two-logit tensor even for third-party architectures that can expose auxiliary
outputs.

\section{Confusion Matrix Convention}

Traversable is always the positive class. The aggregated confusion matrix is
\begin{equation}
\mathbf{C}=
\begin{bmatrix}
 \mathrm{TN} & \mathrm{FP}\\
 \mathrm{FN} & \mathrm{TP}
\end{bmatrix},
\label{eq:confusion_convention}
\end{equation}
Here, $\mathbf{C}$ is the test-wide count matrix; TN, FP, FN, and TP denote
true-negative, false-positive, false-negative, and true-positive pixel counts.
Rows are ground truth and columns are prediction. This ordering is
published because false-safe and false-block terminology can otherwise obscure
which denominator is used.

Figure~\ref{fig:confusion_matrix_visual} reads the same matrix as two
operational questions. The upper ground-truth row contains all pixels that
should be blocked and supplies the denominator of FSR. The lower row contains
all pixels that should be traversable and supplies the denominator of FBR.

\begin{figure}[H]
\centering
\resizebox{\textwidth}{!}{\begin{tikzpicture}[
  font=\sffamily\small,
  cell/.style={draw=black!65,very thick,minimum width=3.05cm,
    minimum height=1.78cm,align=center,inner sep=5pt},
  metric/.style={draw,rounded corners=2pt,very thick,align=center,
    text width=4.05cm,minimum height=1.62cm,inner sep=4pt,
    font=\sffamily\footnotesize},
  link/.style={-{Latex[length=2.2mm]},thick}
]
  \node[font=\bfseries] at (5.05,4.90) {Prediction};
  \node[font=\bfseries,align=center] at (0.85,4.34) {Ground\\truth};
  \node[font=\bfseries] at (3.45,4.30) {Untraversable};
  \node[font=\bfseries] at (6.55,4.30) {Traversable};
  \node[font=\bfseries,anchor=east] at (1.78,3.05) {Untraversable};
  \node[font=\bfseries,anchor=east] at (1.78,0.65) {Traversable};

  \node[cell,fill=good!16] (tn) at (3.45,3.05)
    {\textbf{TN}\\correctly\\blocked};
  \node[cell,fill=safety!15] (fp) at (6.55,3.05)
    {\textbf{FP: false safe}\\blocked terrain\\opened};
  \node[cell,fill=warn!18] (fn) at (3.45,0.65)
    {\textbf{FN: false block}\\valid terrain\\removed};
  \node[cell,fill=pidP!13] (tp) at (6.55,0.65)
    {\textbf{TP}\\correctly\\traversable};

  \node[draw=safety,dashed,rounded corners=2pt,fit=(tn)(fp),inner sep=3pt] (rowu) {};
  \node[draw=warn!90!black,dashed,rounded corners=2pt,fit=(fn)(tp),inner sep=3pt] (rowt) {};
  \node[metric,draw=safety,fill=safety!8] (fsr) at (11.65,3.05)
    {\textbf{False-Safe Rate}\\$\FSR=FP/(FP+TN)$\\
    fraction of truly blocked pixels opened};
  \node[metric,draw=warn!90!black,fill=warn!10] (fbr) at (11.65,0.65)
    {\textbf{False-Block Rate}\\$\FBR=FN/(FN+TP)$\\
    fraction of truly valid pixels removed};
  \draw[link,draw=safety] (rowu.east) -- (fsr.west);
  \draw[link,draw=warn!90!black] (rowt.east) -- (fbr.west);
\end{tikzpicture}
}
\caption[Visual interpretation of the binary confusion matrix.]{Visual
interpretation of the binary confusion matrix and the two asymmetric error
rates. A false safe opens annotated blocked terrain; a false block removes
annotated traversable terrain. Original schematic drawn for this thesis.}
\label{fig:confusion_matrix_visual}
\end{figure}

\begin{table}[H]
\centering
\caption{Operational interpretation of the four binary outcomes.}
\label{tab:confusion_interpretation}
\small
\begin{tabularx}{\textwidth}{@{}lllX@{}}
\toprule
\textbf{Ground truth} & \textbf{Prediction} & \textbf{Count} & \textbf{Interpretation} \\
\midrule
Untraversable & Untraversable & TN & Correctly preserves a blocked pixel \\
Untraversable & Traversable & FP & False safe: opens a pixel outside the annotated capability envelope \\
Traversable & Untraversable & FN & False block: removes an annotated feasible pixel \\
Traversable & Traversable & TP & Correctly preserves a feasible pixel \\
\bottomrule
\end{tabularx}
\end{table}

Table~\ref{tab:confusion_interpretation} states the corresponding operational
meaning of each matrix entry.

All principal test metrics use one matrix accumulated over every pixel in the
split. They are therefore micro-aggregated by pixel. An image with more pixels
would contribute more decisions, although fixed evaluation resizing makes all
544 test images equal in pixel count. Per-image metrics are used only for
qualitative ranking.

\section{Accuracy and Asymmetric Error Metrics}

For class $c$, Intersection over Union is
\begin{equation}
 \operatorname{IoU}_c=
 \frac{\operatorname{TP}_c}
 {\operatorname{TP}_c+\operatorname{FP}_c+\operatorname{FN}_c}.
\end{equation}
Here, $c$ denotes one semantic class, and $\operatorname{TP}_c$,
$\operatorname{FP}_c$, and $\operatorname{FN}_c$ count pixels treated as true
positives, false positives, and false negatives when $c$ is considered the
positive class. The operator $\operatorname{IoU}_c$ is therefore the
intersection of prediction and target for $c$ divided by their union.
For the binary convention this becomes
\begin{align}
 \operatorname{IoU}_{\mathrm{trav}}
 &=\frac{\mathrm{TP}}{\mathrm{TP}+\mathrm{FP}+\mathrm{FN}},\\
 \operatorname{IoU}_{\mathrm{untrav}}
 &=\frac{\mathrm{TN}}{\mathrm{TN}+\mathrm{FP}+\mathrm{FN}},\\
 \operatorname{mIoU}
 &=\frac{\operatorname{IoU}_{\mathrm{trav}}
 +\operatorname{IoU}_{\mathrm{untrav}}}{2}.
\end{align}
The subscripts $\mathrm{trav}$ and $\mathrm{untrav}$ denote traversable and
untraversable. In the latter equation, TN becomes the intersection because the
negative class in the global convention is treated as the class of interest.
The unweighted mean prevents the more frequent class from completely
determining the headline score.

Traversable precision, recall, and F1 are
\begin{align}
 P&=\frac{\mathrm{TP}}{\mathrm{TP}+\mathrm{FP}},&
 R&=\frac{\mathrm{TP}}{\mathrm{TP}+\mathrm{FN}},&
 F1&=\frac{2PR}{P+R}.
\end{align}
Here, $P$ is traversable precision, $R$ is traversable recall, and $F1$ is
their harmonic mean. F1 describes positive-class overlap through precision
and recall but does not
describe true-negative performance. It is therefore reported beside rather
than instead of mIoU.

The two safety-oriented rates are
\begin{align}
 \FSR&=\frac{\mathrm{FP}}{\mathrm{FP}+\mathrm{TN}},
 \label{eq:final_fsr}\\
 \FBR&=\frac{\mathrm{FN}}{\mathrm{FN}+\mathrm{TP}}.
\label{eq:final_fbr}
\end{align}
In these equations, $\FSR$ is False-Safe Rate and $\FBR$ is False-Block Rate;
the four count symbols retain the definitions in
Table~\ref{tab:confusion_interpretation}. FSR conditions on truly
untraversable pixels and asks what fraction the model
opens. FBR conditions on truly traversable pixels and asks what fraction the
model blocks. FBR is equal to one minus traversable recall; it is retained as
a named metric because its operational meaning is clearer in this task.

The two errors matter differently. A connected false-safe area can allow a
planner to consider terrain outside the vehicle's annotated capability, while
a false block can make a feasible route disappear or cause a conservative
stop. Reducing one by moving the threshold usually increases the other. A
useful system must therefore report both. Neither rate includes distance from
the vehicle, trajectory intersection, connected-component size, severity of
terrain, prediction confidence, or actual vehicle outcomes. They are
safety-oriented diagnostic metrics, not safety probabilities.

\section{Multi-Seed Reporting}

For the nine-model comparison, the metric for each model and policy is reported as
mean $\pm$ population standard deviation over three fixed seeds. Population
standard deviation is used because the purpose is to describe those three
executed runs, not to estimate an unknown variance from a random sample.

The seed set controls random initialisation, loader order, and deterministic
augmentation sequences. It does not capture dataset-split uncertainty because
all seeds use the same manifest assignment. It also does not capture
hyperparameter uncertainty. A difference smaller than the observed run-to-run
standard deviation is treated cautiously. In the 15-epoch blue+green matrix,
the top FPN/EfficientNet-B0 and U-Net/EfficientNet-B0 means differ by only
0.02 percentage points. In the convergence-aware follow-up,
FPN/EfficientNet-B0 is higher in all three executed seeds and its mean lead is
0.13 points. This is a consistent observed ordering for the declared suite,
but three selected seeds still do not establish universal superiority or
provide a confidence interval. The matched pairwise comparison shows the same
complete consistency: ROD records higher test mIoU than PIDNet-S in all six of
its policy--seed runs.

\section{Qualitative Review Protocol}

Aggregate metrics reveal how many pixels are wrong but not how those errors
are arranged. The reference PIDNet checkpoint is therefore applied to all 544
test images and each image is assigned:
\begin{itemize}
  \item per-image FSR;
  \item per-image FBR;
  \item combined error $\FSR+\FBR$ for distributional ranking;
  \item boundary-band error.
\end{itemize}

The ground-truth boundary band is formed from the morphological difference
between a two-pixel dilation and a two-pixel erosion. Boundary error is the
fraction of pixels inside that band whose predicted class differs from the
target. The band emphasises local contour placement but can be unstable for
very thin regions and visually ambiguous annotations; it is a diagnostic, not
an alternative headline metric.

The representative-range panel uses mechanically selected images near the
10th, 30th, 50th, 70th, and 90th percentiles of combined per-image error.
Separate ranked panels select the maximum FSR, maximum FBR, and a severe
boundary error. This procedure avoids selecting only visually attractive
successes. The captions report identifiers and rates so the relationship
between the visible mask and the numerical metric can be checked.

The checkpoint is also applied to 130 frames from the external, unlabelled
ALICE mobile-robot sequences (scenario~02, sequences~01--05). Four frames are
displayed to span several path and illumination conditions. ALICE is not used
for training, and no ground-truth masks are available in this study. The panel
can show whether predictions are connected, responsive to scene change, or
visibly implausible; it cannot measure accuracy, FSR, FBR, calibration, or
domain generalisation.

\section{CPU and H100 Timing Protocols}

The CPU platform is an AMD Ryzen~5~5500 with six cores and twelve hardware
threads. PyTorch uses six inference threads, batch size 1, and
$640\times384$ input. Thirty-two test tensors are preloaded. Two complete
warm-up rounds are excluded, followed by five measured repeats, for 160 timed
inferences. Eager execution and \texttt{torch.compile} execution are measured
for all nine comparison models. Compilation occurs during warm-up and is
excluded from the timed samples.

The CPU timer covers model forward, full-resolution logits, softmax, and
thresholded prediction. It excludes image decoding, camera acquisition,
visualisation, and disk output. The reference PIDNet and ROD results are tied
to PyTorch 2.13.0, Python 3.12.12, and their exact checkpoints. ROD retains the
common $640\times384$ external input but resizes it internally to
$1024\times1024$ before its transformer encoder.

The common server measurement uses one NVIDIA H100 NVL 3g.47GB MIG partition,
eager PyTorch FP32, batch size 1, and $640\times384$ external input. The 128
test tensors are preloaded. Five warm-up passes are excluded; ten repeats are
measured with explicit CUDA synchronisation around every forward pass. The
same partition and procedure are used for all nine models. These values
control hardware within the architecture comparison, but an H100 MIG
partition is not an embedded target.

The convergence experiment changes learned parameter values but not the model
graphs, external tensor shape, or prediction operations. Its accuracy values
are therefore paired with the existing architecture-level CPU and H100 timing
measurements rather than rerunning identical inference graphs. Training
wall-clock fields from the convergence histories are excluded: workers were
paused and resumed while sharing the MIG partition, which preserves model and
optimiser state but contaminates elapsed-time measurements.

\section{RTX 5060 TensorRT Protocol}

The reference PIDNet checkpoint and the reconstructed blue+green ROD checkpoint
are exported to the Open Neural Network Exchange (ONNX) format and converted
to fixed-shape, batch-one TensorRT 10.16.1 engines. The builder enables FP16
for eligible operations and retains FP32 for numerically sensitive reduction,
softmax, normalisation, and selected unary operations. The result is described
as FP16-enabled mixed precision rather than all-FP16 inference.

\begin{figure}[H]
\centering
\resizebox{\textwidth}{!}{\begin{tikzpicture}[
  node distance=0.25cm,
  stage/.style={draw,minimum height=0.9cm,text width=1.62cm,align=center,font=\footnotesize,inner sep=2pt},
  included/.style={stage,fill=softgreen},
  excluded/.style={stage,fill=black!8,text=black!55},
  arrow/.style={-{Latex[length=2mm]},thin}
]
\node[excluded] (camera) {camera / disk decode};
\node[included,right=of camera] (resize) {resize and normalise};
\node[included,right=of resize] (h2d) {host-to-device};
\node[included,right=of h2d] (model) {model / TensorRT};
\node[included,right=of model] (post) {softmax and threshold};
\node[included,right=of post] (d2h) {mask to host};
\node[excluded,right=of d2h] (planner) {planning and control};
\draw[arrow] (camera) -- (resize);
\draw[arrow] (resize) -- (h2d);
\draw[arrow] (h2d) -- (model);
\draw[arrow] (model) -- (post);
\draw[arrow] (post) -- (d2h);
\draw[arrow] (d2h) -- (planner);
\draw[very thick,good] ($(resize.south west)+(0,-0.32)$) -- node[below,font=\footnotesize,text=black]{RTX~5060 measured pipeline scope} ($(d2h.south east)+(0,-0.32)$);
\draw[very thick,pidP] ($(model.north west)+(0,0.32)$) -- node[above,font=\footnotesize,text=black]{device/forward scope} ($(post.north east)+(0,0.32)$);
\end{tikzpicture}
}
\caption[Timing scopes used for the RTX 5060 benchmark.]{Timing scopes used for
the RTX~5060 benchmark. Grey stages are excluded. The device/forward scope
starts from GPU-resident input; the measured pipeline starts from an already
decoded host RGB array and returns a binary mask to host memory.}
\label{fig:timing_scopes}
\end{figure}

Figure~\ref{fig:timing_scopes} makes the two reported timing boundaries
explicit before their sample counts are introduced.

Timing uses 64 preloaded native-resolution test frames, ten excluded warm-up
passes, and five measured repeats, yielding 320 samples per model and scope.
The \emph{device scope} includes a device-to-device buffer copy, TensorRT
execution, softmax, and thresholding. It excludes host-to-device input and
device-to-host mask transfer. The \emph{measured pipeline} includes CPU resize
and normalisation, input transfer, TensorRT inference, softmax, thresholding,
and return of the mask to CPU memory. Both exclude disk decode, camera-driver
latency, visualisation, planning, actuation, and one-time engine construction.

After timing, each engine is evaluated on all 544 official test images. This
full accuracy pass is mandatory because mixed precision can move probabilities
across the 0.50 threshold even when average numerical differences are small.
TensorRT metrics are reported separately from the original PyTorch metrics.

\section{Artifact Provenance and Reproduction}

\noindent\textbf{Reproducibility repository:}\\
\url{https://github.com/hax2/cat-pidnet-thesis-reproducibility}. This public
repository is the companion artifact for the thesis. It contains the source
code, exact configurations, split and label-conversion records, experiment
launchers, checkpoint and metric evidence, audit scripts, and instructions
needed to reproduce the reported workflows. The appendix provides direct
commands and expected outputs so that the link is usable without having to
infer the execution order from the source tree.

A headline result is considered identified only when the following artifacts
are named:
\begin{enumerate}
  \item the processed split manifest and binary mapping;
  \item the complete YAML configuration;
  \item the selected \texttt{best.pt} checkpoint;
  \item the evaluator and metric implementation;
  \item the saved JSON result and, where applicable, hardware report.
\end{enumerate}

The reference PIDNet chain includes SHA-256 identities for all five elements and
for the qualitative review, CPU benchmark, ONNX model, TensorRT engine, and
TensorRT report. The saved test confusion matrix was independently recomputed
in a fresh CPU evaluation. Exact equality of the counts establishes
\emph{evaluation reproducibility}: the same artifact and evaluator produce the
same score. It is different from \emph{training repeatability}, which asks how
new stochastic training runs vary and is addressed only by the separate
three-seed studies.

The blue-only fixed-budget architecture ledger contains all 27 configurations,
metric records, and checkpoint hashes. The corresponding fixed-budget
blue+green ledger contains all 27 configuration and metric chains and all
confusion-derived integrity checks pass; 21 locally retained checkpoints have
hashes. The six earlier controlled PIDNet-S and ROD blue+green checkpoints
were not included in that download bundle.

The newer convergence-aware ledger is complete at the identity level: it
contains 27 configurations, histories, train summaries, test records, and
server-computed checkpoint SHA-256 values. All 27 confusion matrices reproduce
their reported metrics and each contains 133,693,440 evaluated pixels. The
checkpoint binaries themselves were not copied locally, so the ledger proves
their recorded server identity but does not permit local reevaluation without
retrieving a selected file.

The convergence campaign also records a software-stack limitation. PIDNet-S,
ROD, DDRNet, and BiSeNet ran with PyTorch 2.7.1 and CUDA 12.6, while the FPN,
U-Net, and SegFormer launcher used PyTorch 2.13.0 and CUDA 13.0 because its
environment contained the required SMP dependency. All used the same H100 MIG
partition, data, configurations, and metric implementation. The cross-model
ranking is usable as a recipe-level comparison, but framework version is not
fully controlled across architecture families.

Appendix~\ref{app:reproduction} gives commands and expected outputs.
Appendix~\ref{app:hashes} records the artifact identities. The repository ledgers
provide the per-seed architecture paths and integrity status.
